\documentclass[11pt]{article}

\usepackage[a4paper,margin=1in]{geometry}
\usepackage[T1]{fontenc}
\usepackage{lmodern}
\usepackage{microtype}
\usepackage{amsmath,amssymb,amsthm,mathtools}
\usepackage{booktabs,tabularx}
\usepackage{enumitem}
\usepackage[numbers,sort&compress]{natbib}
\usepackage[colorlinks=true,linkcolor=blue,citecolor=blue,urlcolor=blue]{hyperref}
\usepackage[nameinlink,noabbrev]{cleveref}

\setlength{\emergencystretch}{3em}

\newtheorem{theorem}{Theorem}[section]
\newtheorem{proposition}[theorem]{Proposition}
\newtheorem{corollary}[theorem]{Corollary}
\newtheorem{lemma}[theorem]{Lemma}
\theoremstyle{definition}
\newtheorem{definition}[theorem]{Definition}
\theoremstyle{remark}
\newtheorem{remark}[theorem]{Remark}

\newcommand{\e}{\mathrm e}
\newcommand{\one}{\mathbf 1}
\newcommand{\Lzero}{\mathrm{L0}}
\newcommand{\Lone}{\mathrm{L1}}
\newcommand{\Ltwo}{\mathrm{L2}}

\title{\textbf{Determinant-Two Liouville Pullback}\\
\large An Exact Shift-Two Frame for Literal Affine M\"obius Pairs
and the Growing-Progression Barrier}
\author{Liang Wang}
\date{July 2026}

\begin{document}
\maketitle

\begin{abstract}
Consider positive affine forms
\[
 D(z)=d+sz,\qquad V(z)=u+az,\qquad su-ad=2,
\]
on one literal TPC generic interval.  On the primitive branches
exported in TPC-93, \(a,s\) are positive, coprime, and odd.  The
integer
\[
 n=sV(z)=aD(z)+2=su+as z
\]
therefore gives an order-preserving bijection from the affine
parameter to one residue class modulo \(as\).  Complete
multiplicativity of the Liouville function gives the exact identity
\[
 \boxed{\;
 \mu(D(z))\mu(V(z))
 =
 \mu^2(D(z))\mu^2(V(z))\lambda(as)
 \lambda(n-2)\lambda(n).\;}
\]
We insert this identity into the full literal TPC-108 coefficient,
retaining its coprimality mask, periodic factor, physical weight,
additive phase, interval origin, and outer key.  Thus the
determinant-two signed problem is exactly a shift-two Liouville
correlation on a \emph{growing arithmetic progression}, with two
quotient-squarefree masks.  The pullback itself has no exponent loss
and preserves every prefix.  It does not turn the block into an
unrestricted correlation, and known fixed-form logarithmic theorems
do not supply the required growing-modulus, ordinary-weight,
all-prefix estimate.  The algebra is \(\Lzero\); its
provenance-preserving attachment is \(\Lone\); no \(\Ltwo\)
cancellation is proved.
\end{abstract}

\section{Literal determinant-two input}

Fix an interval \(I\subset\mathbb Z\).  Assume throughout that
\begin{equation}
 D(z)=d+sz>0,\qquad V(z)=u+az>0
 \quad(z\in I),
\label{eq:positive}
\end{equation}
where \(a,s\in\mathbb N\) and
\begin{equation}
 su-ad=2.
\label{eq:det-two}
\end{equation}
The resolved nonempty branches of TPC-93 additionally satisfy
\begin{equation}
 (a,s)=1,\qquad (as,2)=1,\qquad (d,s)=(u,a)=1
\label{eq:primitive}
\end{equation}
\citep{WangTPC93}.  In particular, \(a\) and \(s\) are odd.
The positivity convention in \eqref{eq:positive} avoids extending
\(\lambda\) to signed integers.  A branch on which a form changes
sign must first be partitioned at its zero; this is an exact finite
operation and is not an arithmetic estimate.

For one TPC-108 key \(\xi\), the literal inner coefficient is
\begin{align}
 A_{\xi,X}(z)
 ={}&
 \mu(D(z))\mu(V(z))
 \one_{(B,V(z))=1}
\notag\\
 &\times
 \chi(\tau+Bz)W_X(\tau+Bz),
\label{eq:literal-A}
\end{align}
and its generic transform is
\begin{equation}
 S_{\xi,X}(\alpha)
 =
 \sum_{z\in I}A_{\xi,X}(z)\e^{-2\pi i\alpha z}.
\label{eq:literal-S}
\end{equation}
All symbols in \eqref{eq:literal-A} are literal: \(B,\tau\), the
fixed-period factor \(\chi\), the physical multiplier \(W_X\), and
the interval are not replaced by averaged surrogates
\citep{WangTPC108}.

\section{The exact shift-two frame}

\begin{theorem}[Determinant-two pullback]
\label{thm:pullback}
Under \eqref{eq:positive}--\eqref{eq:primitive}, put
\begin{equation}
 q=as,\qquad n_0=su=ad+2,\qquad n(z)=n_0+qz.
\label{eq:n-frame}
\end{equation}
Then:
\begin{enumerate}[label=(\roman*),leftmargin=2.2em]
\item
\[
 n(z)=sV(z)=aD(z)+2;
\]
\item \(z\mapsto n(z)\) is an order-preserving bijection from \(I\)
onto
\begin{equation}
 \mathcal N_I
 =
 \{n_0+qz:z\in I\}
 =
 \{n\in n(I):n\equiv0\pmod s,\ n\equiv2\pmod a\};
\label{eq:progression}
\end{equation}
\item on this progression,
\[
 D(z)=\frac{n(z)-2}{a},
 \qquad
 V(z)=\frac{n(z)}s;
\]
\item for every \(z\in I\),
\begin{equation}
 \boxed{
 \mu(D(z))\mu(V(z))
 =
 \lambda(q)\lambda(n(z)-2)\lambda(n(z))
 \mu^2\!\left(\frac{n(z)-2}{a}\right)
 \mu^2\!\left(\frac{n(z)}s\right).}
\label{eq:mobius-pullback}
\end{equation}
\end{enumerate}
\end{theorem}

\begin{proof}
The determinant equation gives
\[
 sV(z)-aD(z)
 =s(u+az)-a(d+sz)=su-ad=2.
\]
This proves (i), and \eqref{eq:n-frame} follows by expansion.
Because \(q>0\), the map is injective and order preserving.  Its
image obeys the two displayed congruences.  Conversely, since
\((a,s)=1\), those congruences select one residue modulo \(as\), and
the interval restriction selects exactly \(n(I)\), proving (ii).
Part (iii) is immediate.

For every positive integer \(m\),
\(\mu(m)=\mu^2(m)\lambda(m)\).  Complete multiplicativity and
\(\lambda(t)^2=1\) give
\begin{align*}
 \lambda(D(z))\lambda(V(z))
 &=
 \lambda(a)\lambda(s)
 \lambda(aD(z))\lambda(sV(z))\\
 &=\lambda(as)\lambda(n(z)-2)\lambda(n(z)).
\end{align*}
Multiplying by the two squarefree indicators proves
\eqref{eq:mobius-pullback}.
\end{proof}

\begin{remark}[Why \(h_0=2\) is used literally]
For determinant \(h_0\), the same calculation yields
\(sV=aD+h_0\) and a shift-\(h_0\) Liouville pair.  The present paper
fixes \(h_0=2\), and \eqref{eq:primitive} then forces the physical
progression modulus \(as\) to be odd.  No average over shifts and no
replacement of \(2\) by a convenient displacement occurs.
\end{remark}

\section{Arithmetic geometry specific to determinant two}

The pullback also gives two exact structural checks that are useful
when later decompositions are proposed.

\begin{proposition}[Common-divisor and parity classification]
\label{prop:gcd-classification}
For every \(z\in I\),
\[
 \gcd(D(z),V(z))\mid2.
\]
Under \((as,2)=1\), the two affine values have the same parity.
Consequently:
\begin{enumerate}[label=(\roman*),leftmargin=2.2em]
\item if they are odd, their gcd is \(1\);
\item if they are even, their gcd is \(2\);
\item on the simultaneous squarefree support, the even--even case is
impossible; hence \(D(z),V(z)\) are odd and coprime, and
\begin{equation}
 \mu(D(z))\mu(V(z))
 =\mu(D(z)V(z)).
\label{eq:quadratic-compression}
\end{equation}
\end{enumerate}
\end{proposition}

\begin{proof}
Every common divisor divides
\(sV(z)-aD(z)=2\).  Since \(a\) and \(s\) are odd, reduction of the
same identity modulo \(2\) gives \(V(z)\equiv D(z)\pmod2\), proving
the first two cases.  If both values are squarefree and odd, they are
coprime and multiplicativity of \(\mu\) gives
\eqref{eq:quadratic-compression}.  If both were squarefree and even,
write \(D=2D'\), \(V=2V'\), with \(D',V'\) odd.  Dividing the
determinant identity by \(2\) would give
\[
 sV'-aD'=1.
\]
The left side is even because \(a,s,D',V'\) are odd, a
contradiction.
\end{proof}

\begin{remark}[A single sign is not a linear sign]
Equation \eqref{eq:quadratic-compression} compresses the active
two-form sign to one M\"obius value of a reducible quadratic
polynomial.  It does not turn it into \(\mu(cz+r)\).  Therefore a
single-linear-form estimate cannot be inserted merely because only
one \(\mu\) symbol remains.
\end{remark}

\begin{proposition}[Ordered sequence isometry]
\label{prop:isometry}
Let \(c_z\) be any complex sequence on \(I\), and put
\(\widetilde c_{n(z)}=c_z\) on \(\mathcal N_I\).  Then, for every
\(1\le p<\infty\),
\[
 \sum_{z\in I}|c_z|^p
 =
 \sum_{n\in\mathcal N_I}|\widetilde c_n|^p,
\]
and the terminal maximum and endpoint-anchored variation satisfy
\begin{align*}
 \max_{T\in I}\left|\sum_{\substack{z\in I\\z\le T}}c_z\right|
 &=
 \max_{M\in\mathcal N_I}
 \left|\sum_{\substack{n\in\mathcal N_I\\n\le M}}
 \widetilde c_n\right|,\\
 |c_{\max I}|+\sum_{z,z+1\in I}|c_{z+1}-c_z|
 &=
 |\widetilde c_{\max\mathcal N_I}|
 +\sum_{\substack{n,n+q\in\mathcal N_I}}
 |\widetilde c_{n+q}-\widetilde c_n|.
\end{align*}
\end{proposition}

\begin{proof}
The map \(z\mapsto n_0+qz\) is an order-preserving bijection, so
each identity is a relabeling of the same finite list.
\end{proof}

The last proposition shows that the sharp signed-prefix/BV norm of
TPC-122 is transported exactly \citep{WangTPC122}.  Measuring
variation between consecutive ambient integers instead of
consecutive progression points would change the norm and is not
authorized by the pullback.

\section{Complete coefficient transport}

Define, for \(n\in\mathcal N_I\),
\begin{align}
 z(n)&=\frac{n-n_0}{q},\notag\\
 \mathcal Q_{a,s}(n)
 &=
 \mu^2\!\left(\frac{n-2}{a}\right)
 \mu^2\!\left(\frac ns\right),\notag\\
 \mathcal W_{\xi,X}(n)
 &=
 \one_{(B,n/s)=1}
 \chi(\tau+Bz(n))W_X(\tau+Bz(n)).
\label{eq:transported-data}
\end{align}

\begin{corollary}[Literal block in the shift-two frame]
\label{cor:literal}
The transform \eqref{eq:literal-S} has the exact representation
\begin{equation}
\boxed{
 S_{\xi,X}(\alpha)
 =
 \lambda(q)
 \sum_{n\in\mathcal N_I}
 \mathcal Q_{a,s}(n)
 \lambda(n-2)\lambda(n)
 \mathcal W_{\xi,X}(n)
 \e^{-2\pi i\alpha z(n)}.}
\label{eq:literal-pullback}
\end{equation}
Every initial \(z\)-prefix corresponds to the same initial
\(n\)-prefix, and
\begin{equation}
 \sum_{z\in I}|A_{\xi,X}(z)|
 =
 \sum_{n\in\mathcal N_I}
 \mathcal Q_{a,s}(n)|\mathcal W_{\xi,X}(n)|.
\label{eq:mass-preserved}
\end{equation}
\end{corollary}

\begin{proof}
Insert \eqref{eq:mobius-pullback} in
\eqref{eq:literal-A}--\eqref{eq:literal-S}.  The mask and all
remaining factors become \eqref{eq:transported-data}.  The map in
\cref{thm:pullback} is an order-preserving bijection, which proves
both the prefix assertion and \eqref{eq:mass-preserved}.
\end{proof}

\begin{remark}[The phase is not silently simplified]
One may write
\[
 \e^{-2\pi i\alpha z(n)}
 =
 \e^{-2\pi i\alpha(n-n_0)/q}
\]
on \(\mathcal N_I\).  The quotient is integral there.  We retain
\(z(n)\) in \eqref{eq:literal-pullback} so that no choice of a
representative for \(\alpha/q\) changes the stated physical phase.
\end{remark}

\section{What the pullback does and does not prove}

\begin{proposition}[No algebraic exponent loss]
\label{prop:no-loss}
Let
\[
 V_{\xi,X}=\sum_{z\in I}|A_{\xi,X}(z)|.
\]
Any estimate proved for the right side of
\eqref{eq:literal-pullback} in its displayed progression, masks,
weight, phase, and all-prefix quantifiers,
\[
 |S_{\xi,X}(\alpha)|
 \le X^{-\eta+o(1)}V_{\xi,X},
\]
returns the same exponent \(\eta\) to the original block.
\end{proposition}

\begin{proof}
The two sums are equal and their absolute comparison masses agree
by \eqref{eq:mass-preserved}.
\end{proof}

\begin{proposition}[Growing-progression firewall]
\label{prop:firewall}
The identity \eqref{eq:literal-pullback} does not by itself reduce
the required estimate to an unrestricted sum
\(\sum_{n\le y}\lambda(n-2)\lambda(n)\).  A valid returning theorem
must retain at least:
\begin{enumerate}[label=(\alph*),leftmargin=2.2em]
\item the actual, generally growing modulus \(q=as\) and residue
\(n_0\bmod q\);
\item both quotient-squarefree masks \(\mathcal Q_{a,s}\);
\item the actual interval origin and every required prefix;
\item the coprimality and periodic masks in \(\mathcal W_{\xi,X}\);
\item the physical weight and additive phase; and
\item uniformity over the predeclared literal outer family.
\end{enumerate}
\end{proposition}

\begin{proof}
Each item occurs explicitly in \eqref{eq:literal-pullback}.  Dropping
one changes either the support or the coefficient.  Equality to an
unrestricted correlation would already fail for the constant
coefficient when \(q>1\), since only one residue class is sampled.
\end{proof}

The logarithmically averaged two-point theorem applies to fixed
affine forms with reciprocal weight and its own terminal-window
quantifiers \citep{Tao2016LogChowla}.  It is not an
ordinary-weight, power-rate theorem uniform in the growing data of
\cref{prop:firewall}.  This observation is a type check, not a
negative statement about that theorem.

\section{A finite determinant-two example}

Take
\[
 a=3,\qquad s=5,\qquad d=u=1.
\]
Then \(su-ad=2\), \(q=15\), and
\[
 D(z)=1+5z,\qquad V(z)=1+3z,\qquad n(z)=5+15z.
\]
The first few values illustrate both exact zeros and surviving
signs:
\begin{center}
\begin{tabular}{@{}rrrrr@{}}
\toprule
\(z\)&\(n\)&\(D\)&\(V\)&\(\mu(D)\mu(V)\)\\
\midrule
0&5 &1 &1 & \(1\)\\
1&20&6 &4 & \(0\)\\
2&35&11&7 & \(1\)\\
3&50&16&10& \(0\)\\
4&65&21&13& \(-1\)\\
\bottomrule
\end{tabular}
\end{center}
For example, at \(z=4\),
\[
 \lambda(15)\lambda(63)\lambda(65)=-1
 =\mu(21)\mu(13),
\]
and both quotient-squarefree masks equal \(1\).  At \(z=1\), the
Liouville product is still a sign, but the quotient-squarefree mask
vanishes because \(V=4\).  This is why that mask cannot be dropped
after the shift-two sign has been exposed.

\section{Stop/go certificate}

For a pulled-back block let \(N=|I|\) and \(q=as\).

\paragraph{GO.}
Continue through this frame only if a later theorem is stated on the
actual progression and controls all literal masks, weights, phases,
origins, and prefixes.  Because \cref{prop:no-loss} is exact, a
proved power saving then returns without an additional algebraic
loss.

\paragraph{SHORT STOP.}
If \(N\) is below the predeclared long-block threshold, the
arithmetic-progression reformulation creates no new averaging.  The
block remains in the short-fiber ledger.

\paragraph{UNIFORMITY STOP.}
If the modulus, mask, or interval range lies outside a cited
theorem's uniformity, that theorem is not callable.  One may not
replace \(q\) by a fixed modulus or the progression by all integers.

\paragraph{ARITHMETIC STOP.}
The statement
\[
 \max_{T\in I}
 \left|
 \sum_{\substack{n\in\mathcal N_I\\z(n)\le T}}
 \mathcal Q_{a,s}(n)\lambda(n-2)\lambda(n)
 \mathcal W_{\xi,X}(n)\e^{-2\pi i\alpha z(n)}
 \right|
 \ll X^{-\eta}V_{\xi,X}
\]
with a fixed \(\eta>0\) is a new fixed-shift arithmetic input.  It is
not proved here and must be labelled \(\Ltwo\) if established.

\section{Claim boundary}

\begin{center}
\small
\begin{tabularx}{\textwidth}{@{}p{0.46\textwidth}X@{}}
\toprule
Statement & Level and status\\
\midrule
Determinant-two \(n\)-frame and identity
\eqref{eq:mobius-pullback}
& Exact finite algebra (\(\Lzero\)).\\
Order, prefix, phase, mask, and mass preservation
& Exact finite algebra (\(\Lzero\)).\\
Attachment to every verified TPC-93/TPC-108 primitive branch
& Literal fixed-\(2\) interface (\(\Lone\)).\\
Uniform growing-progression cancellation
& Not proved; new arithmetic input (\(\Ltwo\) if proved).\\
H3, positive zero-mode exponent, or physical endpoint
& Not proved.\\
\bottomrule
\end{tabularx}
\end{center}

TPC-127 gives a canonical arithmetic coordinate for the signed
determinant-two block.  It is a strict improvement in definition:
the remaining sign is now the literal shift-two Liouville product,
but on a sparse growing progression and behind quotient-squarefree
masks.  The next finite task is to expand those masks with their
actual congruence and tail costs.  No parity breakthrough,
prime-pair lower bound, Hardy--Littlewood asymptotic, or twin-prime
conclusion is claimed.

\begingroup
\footnotesize
\raggedright
\bibliographystyle{plainnat}
\bibliography{references}
\endgroup

\end{document}
